\documentclass[10pt,a4paper,oneside,headings=big,headsepline]{scrreprt}

% Sett margins
%\headsep = 50pt

% Watermark
\usepackage{draftwatermark}
\SetWatermarkText{Internal use only}
\SetWatermarkFontSize{0.15\paperwidth}

% Package to use
\usepackage{fontspec}
\usepackage[english]{babel}
\usepackage{graphicx}
\usepackage{float}
\usepackage{booktabs}
\usepackage{xcolor}
\usepackage[automark,markcase=upper]{scrlayer-scrpage}
\usepackage[hidelinks]{hyperref}

% Basic settings
\author{CTEK E-Mobility AB}
\title{CCU \nanogrid\ Configuration Manual}
\subtitle{Pd\_CM\_004. Version A-4+A-5}

\newcommand{\nanogrid}{NANOGRID\texttrademark}
\newcommand*{\fullref}[1]{\hyperref[{#1}]{\autoref*{#1} \nameref*{#1}}}

% Changing header
\makeatletter
\ihead{\@title\\\@subtitle}
\chead{}
\ohead{
\includegraphics[width=0.2\textwidth]{img/CTEK-E-Mobility.png}}

% Changing footer
\ifoot{}
\cfoot{}
\ofoot{\thepage}
\makeatother

% Add headings and footer on chapter pages.
\renewcommand{\chapterpagestyle}{scrheadings}

\begin{document}
\maketitle
\chapter*{Info}
This document is no longer intended to be shared outside of CTEK E-Mobility
AB.

\tableofcontents


\chapter{Introduction}

This document describes how the load management configuration for a \nanogrid\
master station or CGC is done.

For general information on how to configure the charging station or the CGC
via the web based interface and how to upload configuration files, see the
``CCU Manual (English)''. Appropriate technical knowledge and certification
for work with electrical installations and maintenance work is required to
prevent danger to persons or property.

\section{License}
\label{sec:license}

A valid license for \nanogrid\ is required on all charging stations, as well
as on the \nanogrid\ controller itself.

\section{General architecture}
\label{sec:general-architecture}

A \nanogrid\ site usually consists of
\begin{itemize}
\item A \nanogrid\ controller.
\item One or more charging stations acting as a client to the \nanogrid\
  controller.

\end{itemize}


\chapter{\nanogrid\ Configuration}

\section{Introduction}

\nanogrid\ is a load balancing solution for electrical vehicle (EV)
charging. It will ensure that the available electrical capacity is not
exceeded by the connected \nanogrid\ EVSE units.

\nanogrid\ compliant EVSE units (e.g. CCU R1.95+) are required for a
\nanogrid\ installation. Management of non-\nanogrid\ units in a \nanogrid\
scenario is out of scope for this document.


\section{Configuration}

In order for \nanogrid\ to make correct decisions and actions the
\nanogrid\ installation needs to be configured for \nanogrid.

The most important configuration is the configuration of the \nanogrid\
controller.

The  \nanogrid\ controller unit comes in two variants:

\begin{description}
\item[Grid Controller Unit (GCU)] A unit that is dedicated to the task of load
  balancing. Typically placed in a separate cabinet.
\item[Master Controller Unit (Master)] A unit that is also acts as a charging
  station. The exterior is typically the same as controlled charging stations.
\end{description}

The GCU offers more possible configurations of \nanogrid, as well as is able
to manage more units. For example, the GCU can manage a site with
non-\nanogrid\ units sharing the electrical capacity of the installation, by
either measuring the external load(s) or by measuring the consumption of the
whole site. For more advanced and optional configuration, please refer to
section~\fullref{sec:advanc-conf}.


Documentation contained herein is applicable to both scenarios.

Before the configuration of a \nanogrid\ installation please ensure the
following:

\begin{itemize}
\item \nanogrid\ compliant units are in use
\item Units are equipped with appropriate licenses
\item The local electrical grid conditions are known
\item Electrical installation is finished and will not shortly change
\item Units contain a valid LAN configuration (non-LAN configuration is out of
  scope for this document)
\item Appropriate firmware is in use. As of writing 3.11 or later is
  recommended.
\end{itemize}

\section{GCU/Master configuration -- \texttt{nanogrid.ini}}

The central piece of configuration in a \nanogrid\ installation is the
\nanogrid\ configuration file placed on the \nanogrid\  controller.

This file defines how the nodes are connected electrically as well as
how they may be reached by the local controller unit.

The configuration file is called \verb|nanogrid.ini| in this document and is
in a INI-file format (for more information about the INI-file format see:
\url{https://en.wikipedia.org/wiki/INI\_file}).

Create (or edit an existing) \verb|nanogrid.ini| with a text editor that
understands and respects UNIX-style line-endings. Windows Notepad is NOT
acceptable. For an acceptable editor available for Windows see:
https://notepad-plus-plus.org/ . Do NOT use symbols not covered by ASCII (e.g
åäö).

\subsection{Contents}

The file contains two things:

\begin{itemize}
\item General \nanogrid\ configuration
\item Node configuration
\end{itemize}

General configuration include but is not limited to configuration of the
\nanogrid\ scheduler, see \fullref{sec:advanc-conf-scheduler}.

In scope for this document is two of the most common node types:

\begin{description}
\item [Fuse] A distribution board. The node can not be controlled but imposes
  a limit on the amount of available electrical capacity. (Name \verb|Fuse|
  has historical reasons).
\item[Station] A \nanogrid\ compliant charging station. The node can be directly
  controlled by the \nanogrid\ controller.
\end{description}

\subsection{Distribution boards (Fuse) Format}

A distribution board is described by one of the fuse types. The basic format
of a fuse node is the following:

\begin{verbatim}
[GROUP-1]
type=fuse
rating=32
parent=MAINPANEL
\end{verbatim}

Where

\begin{description}
\item [[GROUP-1]] is the identity of the fuse node. It has to be unique for
  every node in the configuration file.

\item [type=fuse] defines that this node is of the fuse type. The primary
  types covered by this document are \verb|fuse| and \verb|station|. The
  remaining types are covered in section~\ref{sec:feedback-from-non}
  and~\ref{sec:meas-distr-board}.

\item [rating=32] is the upper limit of the node that should not be
  exceeded. This is typically the rating of the fuse (in this case 32
  A). The fuse is assumed to have three phases, all connected up- and
  downstream.

\item [parent=MAINPANEL] defines to which distribution board this station or
  distribution board is connected to (upstream). In this case a distribution
  board with the identity \verb|MAINPANEL|. This means that the node is
  electrically downstream from the node \verb|MAINPANEL|.
\end{description}

\emph{Note:} \verb|station| nodes may not be set as parents as they can only be
end nodes.

Every configuration file describes a tree graph of the electrical
connections that are applicable for the local installation.

This means that every node has a direct parent \emph{except} for one node
which is the grid connection point. The grid connection point is \emph{defined
  by having itself} as a parent in the configuration.

Example:

\begin{verbatim}
[MAINPANEL]
type=fuse
rating=40
parent=MAINPANEL
\end{verbatim}


\subsection{Station format}
\label{sec:station-format}

The basic format of a station node is the following:

\begin{verbatim}
[NG-STATION-1]
type=station
parent=MAINPANEL
outlet/size=2
outlet/1/fallback_current=8
outlet/2/fallback_current=8
PhaseRotation=RST
\end{verbatim}

Where

\begin{description}
\item [[NG-STATION-1]] is the identity of the station node. It has to be
  unique for every node in the configuration file. It also has to correspond
  with the chargeboxidentity parameter set on the charging station.

\item [type=station] defines that this node is of the station type.

\item [parent=MAINPANEL] defines the distribution board this station is
  connected to, in this case a distribution board with the identity
  \verb|MAINPANEL|. This means that the station is electrically downstream from
  the node \verb|MAINPANEL|.

  \emph{Note:} ``station'' nodes may not be set as parents.

\item[outlet/size=2] defines that the station is equipped with two outlets.

\item[outlet/1/fallback\_current=8] integer that states what current should be
  assigned to an EV, in case the charging station loses contact with the
  \nanogrid\ controller. Accepted values are 0 or $>= 6$. The value is in
  Ampere. The fallback current is 8\,A in the example.


\item[outlet/2/fallback\_current=8] integer that states what current should be
  assigned to an EV, in case the charging station loses contact with the
  \nanogrid\ controller. Accepted values are 0 or $>= 6$\. The value is in
  Ampere.  The fallback current is 8\,A in the example.

\item[PhaseRotation=RST] states how the charging station is connected to the
  grid. A combination of R (L1), S (L2), and T (L3). In the example, the
  station is connected L1, L2, and L3.
\end{description}

\emph{Note:} the \verb|outlet/x/fallback_current| parameter \emph{shall} be
set for all EVSE's on the connected charging station.


\subsubsection{Offline clients}
\label{sec:offline-clients}

In the event that a charging station loses contact with the \nanogrid\
controller, it needs to have a specified behaviour. The \nanogrid\ controller
also needs to know how the charging station and it's EVSE's are going to
behave.

The fallback configuration consists of one parameter,
\verb|outlet/X/fallback_current| for a standard EVSE.\footnote{For more complex
charging stations, see section~\ref{sec:fallb-conf}}. Set this paramater to 0
if the EVSE shall not allow charging if it's charging station is
offline. Otherwise, set it to the allowed current allocation, i.e.\ 6 if 6\,A
should be assigned if the charging station goes offline\ldots

\emph{Note:} As of writing it is up to the person configuring the system to
ensure that the combined fallback currents do not assign too much in the event
of a local controller failure.


\subsubsection{Deprecated configuration keys}
\label{sec:depr-conf-keys}

The following keys are no longer being used and can thus be omitted:
\begin{description}

\item [endpoint=mqtt] defines the endpoint of the \nanogrid\ station. The
  endpoint configuration in scope for this document is a mqtt endpoint. In
  this case the node will be accessed via the local MQTT server.

  \emph{Note:} The \verb|endpoint| is not necessary to set from CCU FW r3.6.

\item [outlet/1/connections=1,0,0] defines the first outlet to be a one-phase
  outlet that is electrically connected to L1.

\item [outlet/2/connections=2,3,1] defines the second outlet to be a
  three-phase outlet that has the phase-rotation L2, L3, L1.  Not used since
  the CCU FW r3.11.
\end{description}


\subsection{Installation}

The \verb|nanogrid.ini| must be placed on the local controller unit. And must
ONLY be placed on the local controller unit. Failure to install the file on
the local controller unit or failure to limit installation to the local
controller unit will result in undesired behavior.

The \verb|nanogrid.ini| file may be uploaded via the stations web
interface. Select: ``CCU \nanogrid\ Configuration'' and choose the prepared file
and click ``Upload \nanogrid\ configuration''.  It can also be installed on
the station from a CSMS\footnote{Charging Station Management System -- the
  backend system the Charging Station is connected to.}, by a custom
\verb|DataTransfer.req|. Talk to CTEK E-Mobility AB for more information.

\emph{Note:} If the \verb|nanogrid.ini| file is malformed the firmware will
abort execution.

\section{\nanogrid\ client configuration -- \texttt{chargemanager.ini}}

\emph{Note:} \nanogrid\ node configuration has to take place on every node
that is to be managed by the local controller. The configuration needs to take
place even on a master controller.

Regular configuration of \verb|chargemanager.ini| is out of scope for this
document. It is assumed to have been completed correctly.

From the web interface ensure that \emph{\nanogrid\ Load Balancing} is set to
\emph{Full}. In the \verb|chargemanager.ini| file the configuration key is
\verb|localcontroller| and shall be set to \verb|1|.


\section{Verification}

Once the local controller and all \nanogrid\ stations have been
configured, it is time to verify the installation.

Should a station fail to start a charge, faulty configuration or
network issues may be suspected.

Should the electrical grid have been incorrectly described this will
either result in sub-optimal performance or a risk of tripping fuses.

\section{Complete Example}

Below is a complete example of a configuration file. The configuration
starts with a 63 A main fuse with two 35 A sub-fuses. On each 35 A
sub-fuse two 16 A single phase stations are connected.

\begin{verbatim}
[MAINPANEL]
type=fuse
rating=63
parent=MAINPANEL

[GROUP-1]
type=fuse
rating=35
parent=MAINPANEL

[GROUP-2]
type=fuse
rating=35
parent=MAINPANEL

[STATION1]
type=station
parent=GROUP-1
outlet/size=2
outlet/1/fallback_current=6
outlet/2/fallback_current=0
PhaseRotation=RTS

[STATION2]
type=station
parent=GROUP-1
outlet/size=2
outlet/1/fallback_current=6
outlet/2/fallback_current=0
PhaseRotation=TSR

[STATION3]
type=station
parent=GROUP-2
outlet/size=2
outlet/1/fallback_current=6
outlet/2/fallback_current=0
PhaseRotation=SRT

[STATION4]
type=station
parent=GROUP-2
outlet/size=2
outlet/1/fallback_current=0
outlet/2/fallback_current=0
PhaseRotation=RTS
\end{verbatim}

\section{Advanced Configuration}
\label{sec:advanc-conf}

For some scenarios more advanced \nanogrid\ configuration is required.

\subsection{Scheduler}
\label{sec:advanc-conf-scheduler}

The scheduler controls how scheduling will take place. As of writing
three schedulers exist.

\begin{description}
\item[Equal] Attempts to maintain a ``fair'' distribution
\item[Fifo] First in first out scheduling
\item[SimpleFeedback] A feedback based scheduler, can be abreviated as
  \verb|SFB| in \verb|nanogrid.ini|.
\end{description}

The first two schedulers are designed for open-loop \nanogrid\
scenarios. If the site is equipped with working energy meters the
\verb|SimpleFeedback| scheduler should be prefered.

The scheduler is controlled by setting the \verb|scheduler| parameter at
the top-level category \verb|[General]|.

\subsubsection{Example}

\begin{verbatim}
[General]
scheduler=SIMPLEFEEDBACK

...
\end{verbatim}

\subsection{Phase-switching}

If a station is phase-switching it has the capability to either
provide three-phase current or a single phase which it can choose as
appropriate.

An outlet is configured as phase-switching by receiving that information from
the charging station itself. Previously it was done as in the example below.

\emph{Note:} As of writing not all schedulers support
phase-switching. Only the FIFO scheduler and SIMPLEFEEDBACK scheduler are
supported.

The only configuration needed to be done for a phase-switching station in
\verb|nanogrid.ini| is the extra fallback configuration described in
section~\ref{sec:fallb-conf}.


\subsubsection{Deprecated configuration}
\label{sec:depr-conf-phase-switching}

It is no longer necessary to state that the EVSE has phase-switching
capabilities, as that will be reported by the charging station
itself. However, in versions up to CCU FW r3.11, this was done by:
\begin{description}
\item[outlet/X/connections="switch"] stated that the EVSE had phase switching
  capabilities.
\end{description}
where \verb|X| was 1 or 2.


\subsection{Fallback configuration}
\label{sec:fallb-conf}

In the event that a managed node loses contact with a local controller it will
not be able to charge any EV unless a fallback configuration has been
specified.

The fallback configuration consists of one or two parameters depending on the
type of node.

For non-phase switching nodes \emph{only} a fallback current, see
section~\ref{sec:station-format}, is required while phase switching charging
stations requires a fallback output as well.

\emph{Note:} As of writing it is up to the person configuring the system to
ensure that the combined fallback currents do not assign too much in the event
of a local controller failure.


\subsubsection{Fallback behaviour for phase switching EVES's}
\label{sec:fallb-behav-phase}

The fallback output tells a phase switching station on which output should be
provided when offline. This \emph{only} makes sense for a phase switching
station and will be disregarded for other charging stations. The value has the
following possible values:

\begin{description}
\item[0] No output
\item[1] Three phase output
\item[2] Single phase L1
\item[3] Single phase L2
\item[4] Single phase L3
\end{description}

\subsubsection{Example}

\begin{verbatim}
...

[STATIC-1]
type=station
parent=MAINPANEL
outlet/size=2
outlet/1/fallback_current=8
outlet/2/fallback_current=8
PhaseRotation=RTS

[SWITCHING-A]
type=station
parent=MAINPANEL
outlet/size=2
outlet/1/fallback_current=8
outlet/1/fallback_output=2
outlet/2/fallback_current=8
outlet/2/fallback_output=3
PhaseRotation=TRS

...
\end{verbatim}


\subsection{Feedback from non-EVSE loads}
\label{sec:feedback-from-non}

If a node has an external load that is measured the scheduler can take
it into consideration. For example assume that a fuse both have a
charging station and large fan attached. Then if the load of the fan
is measured, the scheduler can take it into account when controlling the
charging station.

\subsubsection{Example}

\begin{verbatim}
...

[SHARED-BOARD]
type=measuredfuse
meter="modbus/1/1"
rating=20
parent=MAINPANEL

[STATION-1]
type=station
parent=SHARED-BOARD
outlet/size=2
outlet/1/fallback_current=8
outlet/2/fallback_current=8
PhaseRotation=RTS

...
\end{verbatim}


\subsection{Distribution board measuring the whole load}
\label{sec:meas-distr-board}

If a node has an energy meter attached that measures the load at that
point, i.e the aggregated load of the sub nodes the scheduler may use
that information to make its scheduling.

This is in contrast to the measured fuse where only the external load
is measured. In essence the aggregated fuse is the sum of non-EV loads
and EV loads.

\subsubsection{Example}

\begin{verbatim}
...

[SHARED-BOARD]
type=aggregatedfuse
meter="modbus/1/1"
rating=20
parent=MAINPANEL

[STATION-1]
type=station
parent=SHARED-BOARD
outlet/size=2
outlet/1/fallback_current=6
outlet/2/fallback_current=6
PhaseRotation=RTS

...
\end{verbatim}

\subsection{EMS Configuration}

If a node has an EMS (Energy Management System) attached that announces the
available current  in the grid, the scheduler may use that information
to make its scheduling. In this case, the available current through the
distribution board will be set to the \emph{minimum} of the \verb|rating| and
what the EMS reports as available.

EMS devices can be configured on a distribution board, i.e.\ nodes of type:

\begin{itemize}
\item \verb|fuse|
\item \verb|measuredfuse|
\item \verb|aggregatedfuse|
\end{itemize}

An EMS is added with the configuration:
\begin{verbatim}
...

[MAINPANEL]
type=fuse
rating=200
parent=MAINPANEL
ems=ehub:192.168.200.12
emsfallback=100

...
\end{verbatim}

Where

\begin{description}
\item [[MAINPANEL]] is the fuse node to add the EMS to.

\item [ems=ehub:192.168.200.12] defines the endpoint of the EMS.
It consists of the EMS \verb|<type>:<address>|. In the example the type is
\verb|ehub| and the address is \verb|192.168.200.12|.

\item [emsfallback=100] defines the fallback current in ampere
to use if connection is lost with the EMS device.
\end{description}

\emph{Note:} Currently the only supported device is the
Ferroamp EnergyHub, which is enabled by using the device type
\verb|ehub|.

\emph{Note:} The emsfallback configuration is a required parameter.

\emph{Note:} The EMS is considered offline (fallback is used)
if no successful communication with the device has occured during a 30
second period.


\subsection{IT-net}
\label{sec:it-net}


\subsubsection{Electrical connection}
\label{sec:electr-conn}

Stations used in IT-net shall be connected according to the following
schematic, figure~\ref{fig:it-net-example}.

\begin{figure}[htbp]
  \includegraphics{nanogrid/figures/it-net_stations_rotation_schema.pdf}
  \caption{Example schematics for IT-net.  Note that you need to verify where
    the right EVSE (connector) physically takes it phase from, as this depends
    on the charging station model.}
  \label{fig:it-net-example}
\end{figure}

\subsubsection{Main fuse configuration}
\label{sec:main-fuse-conf}

NanoGrid does only support main fuses with the same rating. Configure NanoGrid
with a main fuse setting equal to the fuse value divided by 1.732 (the square
root of three) and round the value downwards. Example:

\begin{quote}
  For fuse rating 63A set the fuse configuration parameters to 63/1.732≈36A.
\end{quote}


\subsubsection{EVSE mapping}
\label{sec:EVSE-mapping}

Each station and EVSE must be configured to reflect how the incoming phases
are connected and how the measured current values are mapped to those
phases. Refer to figure~\ref{fig:it-net-example}. It is critical that all
EVSE's follow a regular pattern. The mapping shall be as in
table~\ref{tab:evse-phase-mappings}. Note, it's important to configure the
phase and EVSE mappings correctly, according to the type of the charging
station.

\begin{table}[htbp]
  \caption{EVSE and phase mappings.}
  \label{tab:evse-phase-mappings}
  \begin{tabular}{lrcr}
    \toprule
    Station number &  EVSE & Map current to & Electrical load \\
    \midrule
    1, 4, 7, \ldots & Left  & L1 & L1 to common L2 \\
                    & Right & L3 & L3 to common L2 \\
    \midrule
    2, 5, 8, \ldots & Left  &  L2 & L2 to common L3 \\
                    & Right &  L1 & L1 to common L3 \\
    \midrule
    3, 6, 9, \ldots & Left  &  L3 &  L3 to common L1 \\
                    & Right &  L2 & L2 to common L1 \\
    \bottomrule
  \end{tabular}
\end{table}

\subsubsection{Example configuration}
\label{sec:example-conf}

An example configuration, where there is a 63A fuse in the incoming
distribution panel, called \verb|MAINPANEL|. However, according to the rule in
section~\ref{sec:main-fuse-conf}, \verb|MAINPANEL| is configured with
36A. ($63 / \sqrt{3} = 36$).

Note the comment in the \verb|nanogrid.ini| file. Comments can be added by
starting a line with e.g.\ a \verb|#| character.

2\textsuperscript{nd} note: whether the two ``phases'' connected to each EVSE
(EV connector), is connected to L1 \& L2 on the incoming terminal, or if they
should be connected to L1 \& L3 depends on the type of the charging
station. Check the manual.

\begin{verbatim}
[MAINPANEL]
type=aggregatedfuse
# The rating is the actual fuse 63A / sqrt(3) = 36A, this is done as this
# is a IT-net.
rating=36
parent=MAINPANEL

[CC2-1]
type=station
parent=MAINPANEL
outlet/size=2
outlet/1/fallback_current=8
outlet/2/fallback_current=8
PhaseRotation=RxT

[CC2-2]
type=station
parent=MAINPANEL
outlet/size=2
outlet/1/fallback_current=8
outlet/2/fallback_current=8
PhaseRotation=SxR

[CC2-3]
type=station
parent=MAINPANEL
outlet/size=2
outlet/1/fallback_current=8
outlet/2/fallback_current=8
PhaseRotation=TxS
\end{verbatim}


The example above assumes that the charging stations in question are of the
type Chargestorm Connected 2. Had it been the previous version, Chargestorm
Connected, the \verb|PhaseRotation| lines should have looked differently.

Had \verb|CC2-1| been of type Chargestorm Connected, the configuration for
that station should have looked like:
\begin{verbatim}
[CC2-1]
type=station
parent=MAINPANEL
outlet/size=2
outlet/1/fallback_current=8
outlet/2/fallback_current=8
PhaseRotation=RTx
\end{verbatim}
Note the last line \verb|PhaseRotation=RTx|. The station shall also be
physically connected as \verb|L1, L3, L2|.

\end{document}


%%%% Local variables:
%%%% mode: latex
%%%% TeX-PDF-mode: t
%%%% TeX-engine: xetex
%%%% TeX-source-correlate-mode: t
%%%% End:

